%! Author = lili
%! Date = 6/9/26

% \chapter{VW4 - Regulation der Genexpression in Bakterien I}

% TODO F 2 Zahlen

E. coli werden hat ca.
4100 Gene.
Natürlich werden
die 4100 Genprodukte in aufeinander abgestimmten
Mengen benötigt (z.B.\ wenig von einem
Regulatorprotein, viel von einem Enzym des
zentralen Kohlenstoffwechsels).

\begin{center}
    Was bedingt unterschiedliche Expressions-Niveaus?
\end{center}
% TODO \defin{wasbedingtexpressionsniveaus}
% Ist antwort Stärken der Promotoren oder RBS? F 7

Es gibt ``starke'' und
``schwache'' Promotoren

Je mehr ein Promotor von der
Consensus-Sequenz abweicht, um so
geringer ist die Wahrscheinlichkeit der
Transkription.

% TODO F 4 mehr sinn

Es gibt auch ``starke'' und
``schwache'' \glsfirst{rbs}

% TODO F 5-6 mehr sinn

Bakterien können sich an viele sich ändernde
Umweltbedingungen anpassen (so werden bei
Hitzestress große Mengen an Hitzestressproteinen
benötigt, bei Normaltemperatur nur sehr geringe)

\section{Ebenen der Genregulation}\label{sec:ebenen-der-genregulation}
\begin{figure}
    % TODO F 8 + 9
    \caption{Verschiedene Ebenen der Regulation}\label{fig:figure15}
\end{figure}

\begin{table}[H]
    \centering
    \begin{tabular}{p{2cm} p{5cm} p{5cm}}
        Ebene & \textbf{Transkription} (DNA zu RNA) & \textbf{Protein} \\
        Bezeichnung & Genexpression & Enzymaktivität \\
        günstig für & nachhaltige
        Regulation & schnelle Regulation \\
        Energie & sparsam & aufwendig \\
    \end{tabular}\label{tab:table3}
\end{table}

% TODO F 9 nochmal genau anschauen

\subsection{Genregulation durch alternative Sigma-Faktoren}\label{subsec:genregulation-durch-alternative-sigma-faktoren}
\defin{sigma_faktor}
\paragraph{Rolle des sigma-Faktors bei der Promotererkennung} % TODO F 10 
% TODO F 11-12
\subsection{Regulation der Transkriptionsinitiation in Bakterien}\label{subsec:regulation-der-transkriptionsinitiation-in-bakterien}
\begin{table}[H]
    \centering
    \begin{tabular}{p{2cm} p{5cm} p{5cm}}
    \end{tabular}\label{tab:table2}
\end{table}
\subsubsection{Am Beispiel der Regulation der Tryptophan-Biosynthese
in E. coli durch TrpR}
% TODO F 13-14, 16-18
% TODO karteikarten

% TODO Was mit F 15????




